\begin{appendices}

\addtocontents{toc}{\protect\setcounter{tocdepth}{2}}

\section*{Allgemeines}\label{Allgemeines}

Der Anhang dient dazu, zusätzliche Informationen, die für das Verständnis des Hauptteils nicht essentiell sind, aufzuführen. Hierzu gehören Abbildungen, Tabellen, \textsf{R}-Code, Simulationsergebnisse, Algorithmen, Datensätze und eventuell Beweise. Es sollte trotzdem im Hauptteil auf die Inhalte des Anhangs verwiesen werden (siehe Anhang \ref{Allgemeines}). Es ist sinnvoll den Anhang thematisch zu gliedern.

\section*{Tabellen}

\begin{table}[H]
    \centering
    \caption{Beispieltabelle Anhang (eigene Darstellung)}
    \begin{tabular}{|l|l|}
    \hline
        Tabellenüberschrift 1 &  Tabellenüberschrift 2\\
        \hline
         Beispielinhalt & Beispielinhalt\\
         \hline
         Beispielinhalt & Beispielinhalt\\
         \hline
    \end{tabular}
    \label{tab:my_label1}
\end{table}

\end{appendices}